\begin{proposition}\label{fam:src:C.3.2}\label{fam:st:C.3.2}
The unique degenerate base point at the actual admissibility boundary is
\(t=1\). The boundary constants
\(c=x+i(3/\pi)\log6\) give
\[
 \ker_{\mathbb R}\Pi=\mathbb R(1,3,-3,-x-3/2).
\]
For rational $x+3/2=m/n$ in lowest terms, the quotient is a
finite holomorphic quotient of $E_i\times\mathbb C^*$, with a covering
of degree four. For irrational \(x+3/2\), the closure of the period subgroup is generated
by the real line \(\mathbb R(0,1)\) and the two periods
\((i,(1-i)/2)\), \((1,0)\). Its Hausdorff quotient is a real manifold
homeomorphic to \(\mathbb R\times(S^1)^2\).
\end{proposition}
\Uses{the hypotheses, notation and definitions stated here; \cref{fam:src:C.1.1}; \cref{fam:src:C.2.3}}

\begin{proof}
\ref{fam:src:C.1.1} gives \(\operatorname{Im}\beta=3/2\) at this value of \(c\).
The vector \(\widehat\gamma+3\widehat u-3\widehat w\) maps to
\((0,x+3/2)\), while \(\widehat\delta\) maps to \((0,1)\).
For $x+3/2=m/n$ in lowest terms, change coordinates by
$(\zeta_1,\zeta_2)\mapsto(\zeta_1,\zeta_2+\zeta_1/2)$.
The vectors $2\widehat u-\widehat\delta$ and
$2\widehat w-\widehat\delta$ then have periods $(2i,0)$ and $(2,0)$.
Bezout's identity applied to $\widehat\gamma+3\widehat u-3\widehat w$
and $\widehat\delta$ gives the period $(0,1/n)$.
These three periods generate a sublattice of index four in the rank-three
image: its two remaining generators are represented by
$(i,1/2)$ and $(1,1/2)$, each of order two modulo this sublattice.
The product identification is the explicit map
$([\zeta_1],[\zeta_2])\mapsto([\zeta_1/2],e^{2\pi in\zeta_2})$
from $\mathbb C^2$ modulo these three periods to $E_i\times\mathbb C^*$. Lattice inclusion gives the claimed degree-four map.
For irrational $x+3/2$, its additive subgroup with $1$ is dense in the
indicated real line. The other two displayed periods are independent
modulo that line. This proves the closure and the quotient description.
By \ref{fam:src:C.2.3} this is the only degenerate torus cover at the admissibility
boundary; every other finite-base period matrix retains real rank four.
We have described the torus cover of the order-four filling, including
the actual quotient and its topology.
\end{proof}